\section{Generalized cocycles of a group acting on labels}
\label{sec:projective_generalized}

\begin{definition}[Generalized coboundary and cocycle]
    \label{def:mpug_generalized_cocycle}
    \lean{GeneralizedThreeCochain,
        GeneralizedFourCochain,
        ActionTensorGauge.coboundary,
        LSymbol.coboundary,
        GeneralizedThreeCochain.coboundary,
        LSymbol.IsGeneralizedCocycle}
    \leanok
    \uses{def:mpug_l_symbols, def:mpug_l_gauge}
    Let $G$ act on a set $\mathsf X$ of labels, and let
    $\chi\colon G\times\mathsf X\to\C^\times$ be label-dependent. Its
    generalized coboundary is
    \begin{align}
        (d\chi)^x_{g,h}
         & =\frac{\chi^{h x}_g\chi^x_h}{\chi^x_{gh}}.
        \label{eq:mpug_generalized_coboundary}
    \end{align}
    For a label-dependent scalar $2$-cochain
    $f\colon\mathsf X\times G^2\to\C^\times$, set
    \begin{align}
        (df)^x_{g,h,k}
         & =\frac{f^x_{g,hk}f^x_{h,k}}
            {f^{k x}_{g,h}f^x_{gh,k}}.
        \label{eq:mpug_generalized_two_coboundary}
    \end{align}
    As recorded in \cite{gap:fbc25_two_cochain_coboundary_index_typo}, the
    final group index in the denominator is $k$, correcting the printed
    formula. For a label-dependent scalar $3$-cochain
    $\omega\colon\mathsf X\times G^3\to\C^\times$, the corresponding next
    coboundary is
    \begin{align}
        (d\omega)^x_{g,h,k,l}
         & =\frac{\omega^{l x}_{g,h,k}\omega^x_{g,hk,l}
                \omega^x_{h,k,l}}
            {\omega^x_{gh,k,l}\omega^x_{g,h,kl}}.
        \label{eq:mpug_generalized_three_coboundary}
    \end{align}
    A function $\Omega\colon\mathsf X\times G^2\to\C^\times$ is a
    generalized $2$-cocycle when
    \begin{align}
        \Omega^x_{g,hk}\Omega^x_{h,k}
         & =\Omega^{k x}_{g,h}\Omega^x_{gh,k}.
        \label{eq:mpug_generalized_cocycle}
    \end{align}
    The differential identities $d\Omega=1$ and $d(df)=1$ follow the group
    cohomology conventions of
    \cite[lines~5914--5922]{FrancoRubio2025MPUGauging}, with the
    three-to-four-cochain coboundary given explicitly by
    (\ref{eq:mpug_generalized_three_coboundary}).
    No transitivity or finiteness assumption is imposed on $\mathsf X$.

    % Source anchor: arXiv:2502.20257, lines 5914--5922 and lemma:G_cocycle.
\end{definition}

\begin{lemma}[Generalized coboundaries square to one]
    \label{thm:mpug_generalized_coboundary_squared}
    \lean{ActionTensorGauge.isGeneralizedCocycle_coboundary,
        GeneralizedThreeCochain.coboundary_coboundary}
    \leanok
    \uses{def:mpug_generalized_cocycle}
    The generalized coboundary of every label-dependent scalar $1$-cochain
    is a generalized $2$-cocycle. With the index correction of
    \cite{gap:fbc25_two_cochain_coboundary_index_typo}, every label-dependent
    scalar $2$-cochain $f$ satisfies the identity of
    \cite[line~5922]{FrancoRubio2025MPUGauging},
    \begin{align}
        d(df) & =1.
        \label{eq:mpug_generalized_coboundary_squared}
    \end{align}
\end{lemma}

\begin{proof}\leanok
    \uses{def:mpug_generalized_cocycle}
    Substituting the coboundary formulas, the action law
    $k(lx)=(kl)x$ and associativity identify the shifted labels and
    group products. Every value of the original cochain then occurs once in
    the numerator and once in the denominator, so all factors cancel.
\end{proof}

\begin{definition}[Positive label-dependent cochains]
    \label{def:mpug_positive_generalized_cochains}
    \lean{PositiveUnits.embedding,
        PositiveUnits.IsPositive,
        ActionTensorGauge.IsPositiveValued,
        LSymbol.IsPositiveValued}
    \leanok
    \uses{def:mpug_l_symbols, def:mpug_l_gauge}
    Let $\iota\colon\mathbb R_{>0}\to\C^\times$ be the natural
    multiplicative embedding. A label-dependent scalar $1$-cochain $\chi$ is
    positive when
    \begin{align}
        \chi^x_g & \in\iota(\mathbb R_{>0})
        \label{eq:mpug_positive_one_cochain}
    \end{align}
    for every $x\in\mathsf X$ and $g\in G$. A label-dependent scalar
    $2$-cochain $\Omega$ is positive when
    \begin{align}
        \Omega^x_{g,h} & \in\iota(\mathbb R_{>0})
        \label{eq:mpug_positive_two_cochain}
    \end{align}
    for every $x\in\mathsf X$ and $g,h\in G$. Thus positivity means
    membership in the embedded subgroup of strictly positive real numbers,
    not positivity of the real part.
\end{definition}

\begin{definition}[Hat operation on label-dependent cochains]
    \label{def:mpug_generalized_hat}
    \lean{ActionTensorGauge.hat,
        LSymbol.hat}
    \leanok
    \uses{def:mpug_l_symbols, def:mpug_l_gauge}
    Following the one-index and two-index cases of the hat operation in
    \cite[line~5912]{FrancoRubio2025MPUGauging}, for label-dependent scalar
    cochains define
    \begin{align}
        \widehat\chi^{\,x}_g
         & =\chi^{\,g x}_{g^{-1}},
        \label{eq:mpug_generalized_hat_one}
    \end{align}
    and
    \begin{align}
        \widehat\Omega^{\,x}_{g,h}
         & =\Omega^{\,(gh)x}_{h^{-1},g^{-1}}.
        \label{eq:mpug_generalized_hat_two}
    \end{align}
\end{definition}

\begin{lemma}[Hat involution and generalized coboundaries]
    \label{thm:mpug_generalized_hat_coboundary}
    \lean{ActionTensorGauge.hat_hat,
        LSymbol.hat_hat,
        LSymbol.hat_coboundary}
    \leanok
    \uses{def:mpug_generalized_cocycle, def:mpug_generalized_hat}
    The label-dependent hat operation is an involution:
    \begin{align}
        \widehat{\widehat\chi} & =\chi,
        \label{eq:mpug_generalized_hat_hat_one}
    \end{align}
    and
    \begin{align}
        \widehat{\widehat\Omega} & =\Omega.
        \label{eq:mpug_generalized_hat_hat_two}
    \end{align}
    It commutes with the generalized coboundary:
    \begin{align}
        \widehat{d\chi} & =d\widehat\chi.
        \label{eq:mpug_generalized_hat_coboundary}
    \end{align}
\end{lemma}

\begin{proof}\leanok
    \uses{def:mpug_generalized_cocycle, def:mpug_generalized_hat}
    Applying (\ref{eq:mpug_generalized_hat_one}) twice and cancelling each
    group element with its inverse gives
    (\ref{eq:mpug_generalized_hat_hat_one}), and applying
    (\ref{eq:mpug_generalized_hat_two}) twice gives
    (\ref{eq:mpug_generalized_hat_hat_two}). For
    (\ref{eq:mpug_generalized_hat_coboundary}), the action law
    $g(h x)=(gh)x$ and $g^{-1}\bigl((gh)x\bigr)=h x$ give
    \begin{align}
        (d\widehat\chi)^x_{g,h}
         & =\frac{\widehat\chi^{\,h x}_g\widehat\chi^{\,x}_h}
            {\widehat\chi^{\,x}_{gh}}
        =\frac{\chi^{(gh)x}_{g^{-1}}\chi^{h x}_{h^{-1}}}
         {\chi^{(gh)x}_{(gh)^{-1}}}
        =(d\chi)^{(gh)x}_{h^{-1},g^{-1}}
        =\widehat{d\chi}^{\,x}_{g,h}.
    \end{align}
    The two sides therefore carry the same three values of $\chi$, reordered
    in $\C^\times$.
\end{proof}

\begin{definition}[Regular-action matrices]
    \label{def:mpug_generalized_regular_matrices}
    \lean{LSymbol.regularMatrix}
    \leanok
    \uses{def:mpug_l_symbols}
    Let $\Omega\colon\mathsf X\times G^2\to\C^\times$ be an $L$-symbol.
    As in \cite[Lemma~\texttt{lemma:G\_cocycle}]{FrancoRubio2025MPUGauging},
    define the regular-action matrix $X^x_g$ on the group algebra $\C G$ by
    \begin{align}
        X^x_g\ket{k}
         & =\Omega^{k^{-1}x}_{g,k}\ket{gk}.
        \label{eq:mpug_generalized_regular_matrix}
    \end{align}
    % Source anchor: arXiv:2502.20257, lines 5938--5942.
\end{definition}

\begin{lemma}[Generalized projective multiplication law]
    \label{thm:mpug_generalized_regular_multiplication}
    \lean{LSymbol.regularMatrix_mul}
    \leanok
    \uses{def:mpug_generalized_cocycle,
        def:mpug_generalized_regular_matrices}
    Let $G$ be finite. If $\Omega$
    satisfies~(\ref{eq:mpug_generalized_cocycle}), then
    \begin{align}
        X^{h x}_gX^x_h
         & =\Omega^x_{g,h}X^x_{gh}.
        \label{eq:mpug_generalized_regular_multiplication}
    \end{align}
    % Source anchor: arXiv:2502.20257, lemma:G_cocycle, lines 5931--5943,
\end{lemma}

\begin{proof}\leanok
    \uses{def:mpug_generalized_cocycle,
        def:mpug_generalized_regular_matrices}
    Applying the left side to $\ket{k}$ yields the coefficient
    \begin{align}
        \Omega^{(hk)^{-1}h x}_{g,hk}\Omega^{k^{-1}x}_{h,k}
         & =\Omega^{k^{-1}x}_{g,hk}\Omega^{k^{-1}x}_{h,k} \\
         & =\Omega^x_{g,h}\Omega^{k^{-1}x}_{gh,k},
    \end{align}
    where the last equality is
    (\ref{eq:mpug_generalized_cocycle}) at the label $k^{-1}x$.
    Formula~(\ref{eq:mpug_generalized_regular_matrix}) now gives
    (\ref{eq:mpug_generalized_regular_multiplication}).
\end{proof}

\begin{lemma}[Finite-group power trivialization]
    \label{thm:mpug_generalized_cocycle_power}
    \lean{LSymbol.regularMatrix_eq_permutation_mul_diagonal,
        LSymbol.regularDeterminant_ne_zero,
        LSymbol.determinantCochain,
        LSymbol.pow_card_eq_coboundary_apply,
        LSymbol.pow_card_eq_coboundary}
    \leanok
    \uses{def:mpug_generalized_cocycle,
        def:mpug_generalized_regular_matrices}
    Let $G$ be finite and let $\Omega$ be a generalized $2$-cocycle. Then
    $\Omega^{|G|}$ is a generalized coboundary. More precisely, with
    \begin{align}
        \chi^x_g
         & =\det X^x_g,
        \label{eq:mpug_generalized_determinant_cochain}
    \end{align}
    one has
    \begin{align}
        \left(\Omega^x_{g,h}\right)^{|G|}
         & =\frac{\det X^{h x}_g\det X^x_h}{\det X^x_{gh}}
        =(d\chi)^x_{g,h}.
        \label{eq:mpug_generalized_power_coboundary}
    \end{align}
    This establishes
    \cite[Lemma~\texttt{lemma:G\_cocycle}]{FrancoRubio2025MPUGauging}; as
    recorded in \cite{gap:fbc25_generalized_cocycle_determinant_typo}, the
    second numerator determinant in
    (\ref{eq:mpug_generalized_power_coboundary}) carries the index $h$,
    correcting the printed formula.

    % Source anchor: arXiv:2502.20257, lemma:G_cocycle, lines 5931--5948.
\end{lemma}

\begin{proof}\leanok
    \uses{def:mpug_generalized_regular_matrices,
        thm:mpug_generalized_regular_multiplication}
    Let $P_g$ be the regular-action permutation matrix defined by
    $P_g\ket{k}=\ket{gk}$, and set
    $D^x_g=\operatorname{diag}(\Omega^{k^{-1}x}_{g,k})_{k\in G}$.
    Then
    \begin{align}
        X^x_g
         & =P_gD^x_g.
    \end{align}
    Since $\det P_g\in\{1,-1\}$ and every diagonal entry of $D^x_g$
    is nonzero, $\det X^x_g\ne0$. Thus
    (\ref{eq:mpug_generalized_determinant_cochain}) takes values in
    $\C^\times$. Taking determinants in
    (\ref{eq:mpug_generalized_regular_multiplication}), using
    multiplicativity of the determinant on the left and
    $\det(cX)=c^{|G|}\det X$ for the $|G|\times|G|$ matrix $X^x_{gh}$ on the
    right, gives
    \begin{align}
        \det X^{h x}_g\det X^x_h
         & =\left(\Omega^x_{g,h}\right)^{|G|}\det X^x_{gh}.
        \label{eq:mpug_generalized_determinant_product}
    \end{align}
    Dividing (\ref{eq:mpug_generalized_determinant_product}) by
    $\det X^x_{gh}$ gives
    (\ref{eq:mpug_generalized_power_coboundary}).
\end{proof}

\begin{lemma}[Positive generalized cocycles are coboundaries]
    \label{thm:mpug_positive_generalized_cocycle}
    \lean{PositiveUnits.positiveRpow,
        PositiveUnits.positiveRpow_inv_natCast_pow,
        ActionTensorGauge.positiveRpow,
        LSymbol.positivePrimitive,
        LSymbol.isPositiveValued_positivePrimitive,
        LSymbol.exists_positive_coboundary}
    \leanok
    \uses{def:mpug_generalized_cocycle,
        def:mpug_positive_generalized_cochains,
        thm:mpug_generalized_cocycle_power}
    Let $G$ be finite and let $G$ act on a set $\mathsf X$. If $\Omega$ is a
    positive generalized $2$-cocycle, then there is a positive
    label-dependent scalar $1$-cochain $\chi_0$ such that
    \begin{align}
        \Omega^x_{g,h}
         & =(d\chi_0)^x_{g,h}
        =\frac{\chi^{\,h x}_{0,g}\chi^x_{0,h}}
         {\chi^x_{0,gh}}
        \label{eq:mpug_positive_coboundary}
    \end{align}
    for all $x\in\mathsf X$ and $g,h\in G$.
    This establishes the first assertion of
    \cite[Corollary at lines~5950--5960]{FrancoRubio2025MPUGauging}; as
    recorded in \cite{gap:fbc25_positive_generalized_cocycle_block_index},
    whereas the source writes $\chi_g$,
    Lemma~\ref{thm:mpug_generalized_cocycle_power} constructs $\chi^x_g$
    through (\ref{eq:mpug_generalized_determinant_cochain}) and the label
    $x$ is therefore retained.
    % Source anchor: arXiv:2502.20257, cor:positive_trivial, lines 5950--5956.
\end{lemma}

\begin{proof}\leanok
    \uses{def:mpug_positive_generalized_cochains,
        thm:mpug_generalized_cocycle_power}
    With $\delta^x_g=\det X^x_g$ denoting the determinant cochain from
    Lemma~\ref{thm:mpug_generalized_cocycle_power}, set
    \begin{align}
        \chi^x_{0,g}
         & =\lvert\delta^x_g\rvert^{1/\lvert G\rvert}.
        \label{eq:mpug_positive_determinant_root}
    \end{align}
    As recorded in
    \cite{gap:fbc25_positive_generalized_cocycle_determinant_sign}, because
    the regular-action permutation can contribute a sign to $\delta^x_g$, its
    norm is taken before the positive $\lvert G\rvert$-th root.
    Since $\Omega$ is positive and the complex norm is multiplicative, taking
    norms in (\ref{eq:mpug_generalized_power_coboundary}) gives
    \begin{align}
        \Omega^{\lvert G\rvert} & =d\lvert\delta\rvert.
        \label{eq:mpug_positive_norm_power}
    \end{align}
    Positive real powers are multiplicative, so
    (\ref{eq:mpug_positive_determinant_root}) and
    (\ref{eq:mpug_positive_norm_power}) give
    \begin{align}
        (d\chi_0)^{\lvert G\rvert}
         & =d\lvert\delta\rvert=\Omega^{\lvert G\rvert},
        \label{eq:mpug_positive_root_power}
    \end{align}
    and the $\lvert G\rvert$-th power is injective on $\mathbb R_{>0}$. Hence
    $d\chi_0=\Omega$.
\end{proof}

\begin{lemma}[Hat-compatible positive primitive]
    \label{thm:mpug_positive_generalized_cocycle_hat}
    \lean{ActionTensorGauge.coboundary_div,
        ActionTensorGauge.coboundary_positiveRpow,
        LSymbol.exists_positive_coboundary_hat_mul_eq_one}
    \leanok
    \uses{def:mpug_generalized_hat,
        def:mpug_positive_generalized_cochains,
        thm:mpug_generalized_hat_coboundary,
        thm:mpug_positive_generalized_cocycle}
    Under the hypotheses of
    Lemma~\ref{thm:mpug_positive_generalized_cocycle}, suppose additionally
    that
    \begin{align}
        \Omega^x_{g,h}\widehat\Omega^{\,x}_{g,h}
         & =1
        \label{eq:mpug_positive_hat_hypothesis}
    \end{align}
    for every $x\in\mathsf X$ and $g,h\in G$. Then, as in the second assertion
    of \cite[Corollary at lines~5950--5960]{FrancoRubio2025MPUGauging}, the
    positive cochain can be chosen so that
    \begin{align}
        \Omega & =d\chi,
        \label{eq:mpug_positive_hat_coboundary}
    \end{align}
    and
    \begin{align}
        \chi^x_g\widehat\chi^{\,x}_g
         & =1.
        \label{eq:mpug_positive_hat_witness_identity}
    \end{align}
    % Source anchor: arXiv:2502.20257, cor:positive_trivial, lines 5955--5960.
\end{lemma}

\begin{proof}\leanok
    \uses{thm:mpug_generalized_hat_coboundary,
        thm:mpug_positive_generalized_cocycle}
    Let $\chi_0$ satisfy $d\chi_0=\Omega$ as in
    Lemma~\ref{thm:mpug_positive_generalized_cocycle}. Define
    \begin{align}
        \chi^x_g
         & =\sqrt{\frac{\chi^x_{0,g}}
                {\widehat\chi^{\,x}_{0,g}}}
        =\sqrt{\frac{\chi^x_{0,g}}
             {\chi^{\,g x}_{0,g^{-1}}}}.
        \label{eq:mpug_positive_hat_witness}
    \end{align}
    All values are positive, so the positive square roots are unambiguous.
    Using (\ref{eq:mpug_generalized_hat_coboundary}) and the hypothesis
    (\ref{eq:mpug_positive_hat_hypothesis}) gives
    \begin{align}
        d\chi
         & =\sqrt{d\chi_0/d\widehat\chi_0}
        =\sqrt{\Omega/\widehat\Omega}
        =\sqrt{\Omega^2}
        =\Omega.
    \end{align}
    The same substitution in
    (\ref{eq:mpug_positive_hat_witness}) gives
    \begin{align}
        \widehat\chi
         & =\sqrt{\widehat\chi_0/\chi_0}
        =\chi^{-1},
    \end{align}
    which proves
    (\ref{eq:mpug_positive_hat_witness_identity}).
\end{proof}
